% Part: first-order-logic
% Chapter: tableaux
% Section: provability-quantifiers

\documentclass[../../../include/open-logic-section]{subfiles}

\begin{document}

\olfileid[pt]{fol}{tab}{qpr}
\olsection{\usetoken{S}{derivability} e os quantificadores}

\begin{explain}
  O teorema da completude também exige que as regras de tableaux
  produzam os fatos sobre~$\Proves$ estabelecidos nesta seção.
\end{explain}

\begin{thm}
\ollabel{thm:strong-generalization} Se $c$ é uma constante que não ocorre
em $\Gamma$ nem em~$!A(x)$ e $\Gamma \Proves !A(c)$, então $\Gamma
\Proves \lforall[x][!A(x)]$.
\end{thm}

\begin{proof}
Suponha que $\Gamma \Proves !A(c)$, isto é, existem $!B_1$, \dots, $!B_n
\in \Gamma$ e um !!{tableau} fechado para
\begin{align*}
\{ \sFmla{\False}{!A(c)}, &
\sFmla{\True}{!B_1}, \dots, \sFmla{\True}{!B_n} \}.
\intertext{Devemos mostrar que também existe um !!{tableau} fechado para}
\{ \sFmla{\False}{\lforall[x][!A(x)]}, &
\sFmla{\True}{!B_1}, \dots, \sFmla{\True}{!B_n} \}.
\end{align*}
Considere o !!{tableau} fechado, substitua a primeira hipótese por
$\sFmla{\False}{\lforall[x][!A(x)]}$ e insira
$\sFmla{\False}{!A(c)}$ depois das hipóteses.
\begin{center}
\begin{tableau}{not line numbering}
  [\sFmla{\False}{\formula{A}(c)}
    [\sFmla{\True}{\formula{B}_1}
      [\vdots
      [\sFmla{\True}{\formula{B}_n}, 
        [][][]]]]]
\end{tableau}\qquad
\begin{tableau}{not line numbering}
  [\sFmla{\False}{\lforall[x][\formula{A}(x)]}
    [\sFmla{\True}{\formula{B}_1}
      [\vdots
        [\sFmla{\True}{\formula{B}_n},
          [\sFmla{\False}{\formula{A}(c)}
        [][][]]]]]]
\end{tableau}
\end{center}
O tableau continua fechado, pois todas as !!{sentence}s antes
disponíveis como hipóteses continuam disponíveis no topo do
!!{tableau}. A linha inserida resulta de uma aplicação correta de
$\TRule{\False}{\lforall}$, pois a !!{constant}~$c$ não ocorre em
$!B_1$, \dots, $!B_n$ nem em $\lforall[x][!A(x)]$, isto é, não ocorre
acima da linha inserida no novo !!{tableau}.
\end{proof}

\begin{prop}
\ollabel{prop:provability-quantifiers}
\begin{tagenumerate}{prvEx,prvAll}
\tagitem{prvEx}{$!A(t) \Proves \lexists[x][!A(x)]$.}{}
\tagitem{prvAll}{$\lforall[x][!A(x)] \Proves !A(t)$.}{}
\end{tagenumerate}
\end{prop}

\begin{proof}
\begin{tagenumerate}{prvEx,prvAll}
\tagitem{prvEx}{Um !!{tableau} fechado para
  $\sFmla{\False}{\lexists[x][!A(x)]}, \sFmla{\True}{!A(t)}$ é:
  \begin{oltableau}
    [\sFmla{\False}{\lexists[x][\formula{A}(x)]}, just = \TAss
      [\sFmla{\True}{\formula{A}(t)}, just = \TAss
        [\sFmla{\False}{\formula{A}(t)}, just = {\TRule{\False}{\lexists}[1]}, close]]]
    \end{oltableau}}{}
\tagitem{prvAll}{Um !!{tableau} fechado para
  $\sFmla{\False}{\formula{A}(t)}, \sFmla{\True}{\lforall[x][!A(x)]}, $ é:
  \begin{oltableau}
    [\sFmla{\False}{\formula{A}(t)}, just = \TAss
      [\sFmla{\True}{\lforall[x][\formula{A}(x)]}, just = \TAss
        [\sFmla{\True}{\formula{A}(t)}, just = {\TRule{\True}{\lforall}[2]}, close]]]
    \end{oltableau}}{}
\end{tagenumerate}
\end{proof}

\end{document}
